% Chapter 7 (results half) -- Completion Evaluation: §7.3-§7.7
% Standalone, compilable draft (FPT-template article class + booktabs + float).
% Folded into the thesis main document later via \input.
%
% MERGE / NUMBERING NOTE. This file is the RESULTS half of Chapter 7. Its
% \section{Completion Evaluation} line exists only so the file compiles on its
% own; at final assembly the assembler DROPS this duplicate \section (the chapter
% is opened by sec_7_2_donor_metric.tex, which owns \label{ch:completion-eval}),
% and §7.3-§7.7 below fold in after §7.1/§7.2. Because separate standalone builds
% do NOT share a label namespace, cross-FILE references use typed numerals
% ("Section~\ref{sec:donor-metric}", "Chapter~\ref{ch:completion-method}") rather than \ref{}; only within-file refs use \ref.
%
% Scope: the completion RESULTS. §7.3 surface coverage (donor metric, #32);
% §7.4 downstream box utility + length geometry (#29/#35/#36, #45 cross-ref);
% §7.5 moving cars (#44); §7.6 pre-registered held-out replication on seq 00
% (#42); §7.7 summary + limitations. The metric itself is defined in §7.2; the
% method in Chapter 6.
%
% Honesty rules (personal_agenda P1): per-car medians / paired Wilcoxon; static
% cars by construction; trade-offs stated plainly (per-car length estimate costs
% donor coverage); negatives (#45) as findings, not apologies; measured/inferred
% marked. Numbers traceable to docs/findings.md #29/#32/#35/#36/#37/#41/#42/#44/#45,
% docs/completion/donor_metric.md, and docs/plans/t9{b,c}_*heldout_seq00.md.
%
% Figures (via \graphicspath into the gitignored output/ tree; regenerable):
%   Fig 7.1 donor_metric_08.png              (donor_metric_viz.py)
%   Fig 7.2 completion_box_overlays_08.png   (completion_box_eval_viz.py)
%   Fig 7.3 t11_mover_completions_bev.png    (t11_mover_plausibility.py)
% Floats set [htbp].
%
Section~\ref{sec:donor-metric} established a metric that measures, on real cars, how much
previously-unseen surface a completion recovers, and showed that it passes an
explicit validation battery. This part of the chapter uses it. It reports three
things the completion method was built to deliver and one it was not: that the
completion adds genuinely unseen surface (Section~\ref{sec:results-coverage}),
that this translates into a better amodal bounding box (Section~\ref{sec:results-box}),
that moving cars --- which the static-car metrics cannot score --- complete no
worse than static ones (Section~\ref{sec:results-movers}), and that both headline
claims survive a pre-registered replication on a held-out sequence
(Section~\ref{sec:results-heldout}). Section~\ref{sec:results-summary} draws the
limitations together. Throughout, the unit of analysis is the car: because the
frames of one parked car are autocorrelated, each car is reduced to its per-car
median and significance is tested across cars with the Wilcoxon signed-rank test,
the paired design of Section~\ref{sec:donor-metric}.

\section{Real-data surface coverage}
\label{sec:results-coverage}

The primary completion result is that PCN reconstructs surface the input frame
never observed. Table~\ref{tab:donor-coverage} reports the donor metric on \num{39} of seq~08's
\num{40} well-observed static cars (one car's pairs are all gate-skipped), scoring
three clouds --- the raw partial, its symmetry mirror, and the completion ---
against each car's
donor-only novel set. This is the metric's acceptance run (Finding~\#32), taken
on the production path \emph{before} the Chapter~\ref{ch:completion-method} length prior; the length
geometry is evaluated in Section~\ref{sec:results-box}.

\begin{table}[H]
  \centering
  \caption{Real-data surface coverage on the donor metric (seq~08, per-car
    medians, $n=39$ of the \num{40} well-observed static cars --- one car's pairs
    are all gate-skipped; $\tau=\SI{0.15}{\meter}$;
    Finding~\#32). \texttt{cov@0.1} is the fraction of never-observed (donor-only)
    surface recovered within \SI{0.1}{\meter}; median novel distance is its
    threshold-free companion; out-of-box is the hallucination guard. All three
    pairwise contrasts are significant at $p<10^{-6}$.}
  \label{tab:donor-coverage}
  \begin{tabular}{lccc}
    \toprule
    Cloud & \texttt{cov@0.1} $\uparrow$ & med.\ novel dist.\ (m) $\downarrow$ & out-of-box $\downarrow$ \\
    \midrule
    raw partial       & 0.000          & 0.518          & 0.000 \\
    mirrored          & 0.043          & 0.332          & 0.008 \\
    \textbf{completed} & \textbf{0.304} & \textbf{0.161} & \textbf{0.000} \\
    \bottomrule
  \end{tabular}
\end{table}

The completion recovers about \SI{30}{\percent} of never-observed surface within
\SI{0.1}{\meter} --- \emph{seven times} the symmetry-mirror baseline --- while its
out-of-box fraction stays at the raw partial's floor, so the coverage is real
surface, not surface invented everywhere. The mirror is the informative
comparison: reflecting the input across the car's symmetry plane is a real,
zero-parameter way to add surface, and the learned model earns its place only by
beating it, which it does more than sixfold. This is the project's first positive
real-data evidence that completion adds unseen structure, and it complements the
box-level value of Section~\ref{sec:results-box}.

Two breakdowns qualify the headline. First, coverage is uneven over the car
(Finding~\#32): the occluded \emph{far side} is well recovered (completed cov
\num{0.321}) and the roof reasonably so (\num{0.203}), but the \emph{far end} ---
the portion of the car truncated by the single viewpoint --- is the weakest
region (\num{0.133}). The far end is exactly where a one-sided partial is
normalised around its observed near portion, and it is what the length geometry
of Section~\ref{sec:results-box} sets out to move. (The mirror scores essentially
zero on the far end by construction --- a reflection across the length plane
cannot add end surface --- which is a built-in sanity check on the region split.)
Second, the coverage is robust to input sparsity: completed \texttt{cov@0.1}
holds at \num{0.30}--\num{0.33} from the sparsest (${<}100$ points) to the densest
(${\geq}300$) inputs, while the mirror falls from \num{0.083} on the densest
inputs to \num{0.010} on the sparsest, as the reflection of a sparse cluster
leaves ever more of the car uncovered.
Completion helps most, relative to the cheap baseline, exactly where the sensor
saw least.

\begin{figure}[htbp]
  \centering
  \includegraphics[width=\linewidth]{donor_metric_08.png}
  \caption{The donor metric in bird's-eye view on six seq-08 cars, ordered
    best to worst completed coverage. Grey is the single-frame partial input; the
    coloured points are the donor-only novel set (surface the input frame never
    saw), coloured by whether the completion (green overlay) recovers them within
    \SI{0.1}{\meter}. The best panels close the occluded far side almost entirely;
    the worst isolate the two residual failure modes --- far-end under-completion
    and heading/centre error on diagonal or sparse views --- that
    Section~\ref{sec:results-box} and Chapter~\ref{ch:completion-method} address. Rendered by
    \texttt{scratchpad/donor\_metric\_viz.py}.}
  \label{fig:donor-coverage}
\end{figure}

\section{Downstream utility: the amodal bounding box}
\label{sec:results-box}

Surface coverage shows the completion is geometrically faithful; the box
evaluation shows it is \emph{useful}. For each true-positive (frame, car) pair on
the well-observed static cars we fit an oriented box to the raw partial and to the
completion with the \emph{same} fitter used to build the amodal ground-truth
boxes (so fitter bias cancels), and compare each against the GT box. This is the
pre-registered ``completion adds value'' test of Finding~\#29.
Table~\ref{tab:box-quality} reports it.

\begin{table}[H]
  \centering
  \caption{Amodal box quality, raw partial vs.\ completion (seq~08, per-car
    medians, $n=39$ of the \num{40} well-observed static cars --- one has no
    completed pair; Finding~\#29). BEV IoU is the oriented
    bird's-eye-view box overlap; $|\Delta W|,|\Delta H|,|\Delta L|$ are absolute
    dimension errors; centre error is in the ground plane. $p$ is the Wilcoxon
    signed-rank across cars. These figures use the original Finding~\#29 detection
    configuration; the width, height, and overlap wins reproduce under the promoted
    configuration of Table~\ref{tab:length-geometry} (BEV IoU
    \num{0.725}$\to$\num{0.747}, $|\Delta W|$ \num{0.210}$\to$\num{0.122},
    $|\Delta H|$ \num{0.227}$\to$\num{0.144}).}
  \label{tab:box-quality}
  \begin{tabular}{lcccl}
    \toprule
    Metric & Raw & Completed & Wilcoxon $p$ & \\
    \midrule
    BEV IoU $\uparrow$        & 0.707 & \textbf{0.747} & \num{1.9e-3}  & better \\
    $|\Delta W|$ (m) $\downarrow$ & 0.270 & \textbf{0.170} & \num{1.5e-4}  & better \\
    $|\Delta H|$ (m) $\downarrow$ & 0.255 & \textbf{0.131} & \num{1.6e-10} & better \\
    centre err.\ (m) $\downarrow$ & 0.286 & \textbf{0.234} & \num{2.8e-5}  & better \\
    $|\Delta L|$ (m) $\downarrow$ & 0.447 & 0.456          & 0.65          & neutral \\
    yaw err.\ (deg) $\downarrow$  & 3.5   & 3.0            & 0.74          & neutral \\
    \bottomrule
  \end{tabular}
\end{table}

Completion improves the box on width, height, centre, and overall BEV overlap
with strong significance, and nothing degrades significantly. The improvements
are largest where the sensor saw least --- BEV IoU by input size rises
\num{0.461}$\to$\num{0.599} on the sparsest cars (${<}100$ points),
\num{0.585}$\to$\num{0.678} at \num{100}--\num{300}, and \num{0.703}$\to$\num{0.744}
above \num{300} --- which is the right profile for the use case, since a
well-observed car needs little help. The width gain is not an artefact of the
GT accumulation: it holds both for the \num{13} cars whose width the accumulation
genuinely constrains (both flanks eventually seen: \num{0.291}$\to$\num{0.182})
and for the other \num{26} (\num{0.167}$\to$\num{0.118}).

\paragraph{Length was the exception, and the geometry work fixed it.} In
Table~\ref{tab:box-quality} length is neutral --- the completion neither helps nor
hurts $|\Delta L|$. That was too generous a reading of an untreated defect: the
completion of a partial truncated at the occluded far end stops short, so it
\emph{under}-completes length. Table~\ref{tab:length-geometry} isolates this on
the promoted configuration ($n=40$) and tracks the two Chapter~\ref{ch:completion-method} fixes ---
the fixed \SI{4.14}{\meter} prior (Finding~\#35) and the per-car length estimate
(Finding~\#36) --- against the raw partial and the completion \emph{without} any
length correction.

\begin{table}[H]
  \centering
  \caption{Length geometry on the amodal box (seq~08, well-observed static cars,
    per-car medians, $n=40$, promoted configuration; Findings~\#35/\#36).
    ``no prior'' is the completion with the canonical-frame centre but no
    longitudinal push; ``fixed \SI{4.14}{}'' is the population-median prior;
    ``per-car'' is the final \nth{90}-percentile track estimate. The compact
    row is signed $\Delta L$ on cars ${<}\SI{3.6}{\meter}$ ($n=8$), where a fixed
    prior over-extends.}
  \label{tab:length-geometry}
  \begin{tabular}{lcccc}
    \toprule
    Metric & Raw & no prior & fixed \SI{4.14}{} & \textbf{per-car (final)} \\
    \midrule
    $|\Delta L|$ (m) $\downarrow$   & 0.428 & 0.476 & 0.354 & \textbf{0.304} \\
    BEV IoU $\uparrow$              & 0.725 & 0.743 & 0.747 & \textbf{0.771} \\
    centre err.\ (m) $\downarrow$   & 0.271 & 0.227 & 0.229 & \textbf{0.184} \\
    \midrule
    compact signed $\Delta L$ (m)   & $-0.007$ & --- & $+0.295$ & \textbf{$+0.063$} \\
    \bottomrule
  \end{tabular}
\end{table}

Two things are visible here, and both are stated as the honest picture rather
than the flattering one. First, Table~\ref{tab:box-quality}'s neutral
$|\Delta L|$ averages over- and under-completion and so hides a real regression:
the \emph{signed} length shows the completion systematically \emph{under}-extends
the occluded far end (normal cars $-0.485\to-0.545$; Finding~\#29). On the
promoted configuration this surfaces even in the absolute $|\Delta L|$, where
completion \emph{without} a length prior makes box length \emph{worse} than the
raw partial (\num{0.428}$\to$\num{0.476}, $p=0.027$), which the far-end truncation
causes. The fixed prior
reverses it (\num{0.354}) and the per-car estimate improves it further
(\num{0.304}) while lifting BEV IoU to \num{0.771} and cutting centre error to
\num{0.184}. Second, moving from the fixed prior to the per-car estimate is what
removes the over-extension of genuinely compact cars: their signed $\Delta L$ goes
from \num{+0.295} (the fixed median over-extends a short car) to \num{+0.063}
($p=0.016$ at $n=8$), with the same fix showing up in the per-band hallucination
guard of Section~\ref{sec:donor-metric} (Finding~\#37).

\paragraph{The trade-off, stated plainly.} The per-car estimate buys this box
accuracy at a measured cost in donor coverage: pooled completed \texttt{cov@0.1}
falls \num{0.403}$\to$\num{0.364} and far-end coverage \num{0.346}$\to$\num{0.316}
relative to the fixed prior (Finding~\#36). Per band it is better on compact and
long cars and worse on normal ones. The two metrics genuinely disagree here, and
a control run settled which reading is right: deliberately over-extending
(\nth{90} percentile $+\SI{0.45}{\meter}$) improved \emph{both} coverage and the
box on normal and long cars, which means those completions are still honestly too
short rather than the metric being fooled. In other words the per-car estimate
trades a little pooled coverage --- driven by the normal band it slightly
shortens --- for a correctly-sized box across all bands, and the residual
shortfall is a real limitation, addressed next.

\paragraph{A bounded negative on the residual shortfall.} Targeting each car's
\emph{true} length still leaves completions short on normal and long cars, which
pointed to a second under-extension mechanism inside the completion itself (PCN
under-fills its normalised frame, and the centre push doubles as an output
rescale). Chapter~\ref{ch:completion-method} investigated this under pre-registration (Finding~\#45, the
``closed negative''): decoupling the scale from the length push and adding a
calibrated fill factor did fix the target (normal-car $|\Delta L|$
\num{0.329}$\to$\num{0.232}) but over-extended compacts and failed the
pre-registered compact non-regression guard on both test sequences. The verdict
was do-not-ship; the knobs stay in the code defaulted off. It is reported here
because it \emph{bounds} the result: the residual far-end shortfall on long cars
is a characterised, length-dependent limit, not an oversight.

\begin{figure}[htbp]
  \centering
  \includegraphics[width=\linewidth]{completion_box_overlays_08.png}
  \caption{Amodal box overlays on six seq-08 cars: the ground-truth box (black),
    the box fitted to the raw single-frame partial (blue), and the box fitted to
    the completion (green), over the raw points (grey). The completion tightens
    the box on the occluded width and far end and centres it better; the one
    strong regression panel is a heading error on a sparse (\num{109}-point)
    input, a diagonal-view failure mode also visible in
    Figure~\ref{fig:donor-coverage}. Rendered by
    \texttt{scratchpad/completion\_box\_eval\_viz.py}.}
  \label{fig:box-overlays}
\end{figure}

\section{Moving cars}
\label{sec:results-movers}

Both metrics so far are defined on \emph{static} cars by construction: a paired
amodal reference needs a stationary target. Production nonetheless completes
moving cars, and those completions were never validated. Because completion runs
on a single reference frame's cluster --- not on accumulated points --- there is
no motion smear in its input, so there is a mechanistic reason to expect motion
not to matter; this section checks it (Finding~\#44). On the final seq-08 output
(\num{518} completed car tracks) each track is labelled static, moving, or
ambiguous by the net horizontal displacement of its centroid (a label-free proxy,
since the output tracks are unlabelled), and each completion is scored by a
size-plausibility recipe (bird's-eye dimensions falling in the car range).
Table~\ref{tab:movers} gives the result.

\begin{table}[H]
  \centering
  \caption{Completion plausibility by track motion (final seq-08 output,
    \num{518} completed car tracks; Finding~\#44). Motion is the net displacement
    of the track centroid (median of first vs.\ last five frames). Plausible =
    bird's-eye dimensions in the car range ($L\in[3.3,4.9]$, $W\in[1.5,2.1]$,
    $H\in[1.1,1.7]$~m).}
  \label{tab:movers}
  \begin{tabular}{lcccc}
    \toprule
    Group (net displacement) & $n$ & plausible & rate & median $L/W/H$ (m) \\
    \midrule
    static ($\leq\SI{2}{\meter}$)   & 419 & 225 & \SI{53.7}{\percent} & 3.81 / 1.94 / 1.49 \\
    \textbf{moving ($\geq\SI{5}{\meter}$)} & \textbf{19} & \textbf{11} & \textbf{\SI{57.9}{\percent}} & 3.84 / 2.05 / 1.44 \\
    ambiguous (\SIrange{2}{5}{\meter}) & 80  & 46  & \SI{57.5}{\percent} & 3.95 / 1.92 / 1.49 \\
    all                             & 518 & 282 & \SI{54.4}{\percent} & 3.83 / 1.94 / 1.49 \\
    \bottomrule
  \end{tabular}
\end{table}

Moving cars complete \emph{no worse} than static ones --- slightly higher, in
fact (\SI{57.9}{} vs.\ \SI{53.7}{\percent}), well within noise at $n=19$ --- with
near-identical median dimensions, and every motion bucket clusters at
\SIrange{54}{58}{\percent}. Motion does not degrade completion plausibility, as
the single-frame input predicts. One honest caveat about the absolute level, not
the comparison: the plausibility recipe uses axis-aligned dimensions, which
overestimate width and length for cars oriented diagonally to the world axes, so
the \raisebox{0pt}[0pt][0pt]{$\sim$}\SI{54}{\percent} rate \emph{understates}
completion quality (most failures are inflated width on genuine car-shaped
footprints). Since both groups use the identical recipe, the static-vs-moving
conclusion is unaffected. The value of the check is that the statics-only design
of the earlier sections does not hide a mover failure mode.

\begin{figure}[htbp]
  \centering
  \includegraphics[width=\linewidth]{t11_mover_completions_bev.png}
  \caption{Bird's-eye completions of moving-car tracks on seq~08 (blue partial
    input, green completion). Most form recognisable car footprints; the failures
    are width inflation on diagonally-oriented cars, an artefact of the
    axis-aligned plausibility measure rather than of motion. Rendered by
    \texttt{scratchpad/t11\_mover\_plausibility.py}.}
  \label{fig:movers}
\end{figure}

\section{A pre-registered held-out replication}
\label{sec:results-heldout}

The completion constants --- the width and length priors, the scale correction,
the gate thresholds --- were all developed on seq~08. A fair test of whether the
method \emph{generalises} therefore has to run on a sequence that took no part in
that development, under a protocol fixed in advance so the outcome cannot be
reinterpreted after the fact. We pre-registered one (\texttt{preregistration\_heldout.md},
committed with an immutable timestamp before any held-out run). The held-out
sequence is seq~00 (the intended default, seq~05, had too few well-observed
static cars to test); the two \emph{primary} metrics were fixed as BEV box IoU
(Section~\ref{sec:results-box}) and donor \texttt{cov@0.1}
(Section~\ref{sec:results-coverage}), each requiring the completion to beat the
raw partial with $p<0.05$; three refutation rules (fails to generalise / seq-08
specific length constants / an uncovered result) and a three-level outcome
taxonomy were registered alongside. Nothing was retuned for seq~00.

One confound is logged rather than hidden: seq~00 is in the classifier's training
split, so its \emph{detection} recall reads optimistically. This does not touch
the completion claim, which is a \emph{paired} raw-vs-completed comparison on
identical true-positive inputs, and PCN was trained only on synthetic data and
never saw seq~00. Table~\ref{tab:heldout} gives the outcome on the \num{45}
well-observed static cars.

\begin{table}[H]
  \centering
  \caption{Held-out replication on seq~00 (well-observed static cars, per-car
    medians, $n=45$; the two pre-registered primaries in bold, secondary
    dimensions below; Finding~\#42). Both primaries are decisive wins over the
    raw partial.}
  \label{tab:heldout}
  \begin{tabular}{lccl}
    \toprule
    Metric & Raw & Completed & Wilcoxon $p$ \\
    \midrule
    \textbf{BEV IoU} $\uparrow$          & \textbf{0.739} & \textbf{0.766} & \textbf{\num{1.6e-3}} \\
    \textbf{donor \texttt{cov@0.1}} $\uparrow$ & \textbf{0.000} & \textbf{0.413} & \textbf{${\approx}0$} \\
    $|\Delta W|$ (m) $\downarrow$        & 0.203 & 0.165 & \num{4.0e-4} \\
    $|\Delta H|$ (m) $\downarrow$        & 0.278 & 0.120 & \num{7.8e-9} \\
    \bottomrule
  \end{tabular}
\end{table}

Both headline claims replicate: on a sequence that played no part in tuning, the
completion beats the raw partial on box overlap (\num{0.739}$\to$\num{0.766},
$p=\num{1.6e-3}$) and on surface coverage (\num{0.000}$\to$\num{0.413},
$p\approx0$), with the secondary width and height errors improving as they did on
seq~08. Applying the pre-registered rules verbatim, the arbitration verdict was
\textbf{partially holds} (Finding~\#42): neither generalisation rule triggered ---
both primaries won and no testable length band regressed against its seq-08 level
--- but seq~00's well-observed set contains \emph{no} long cars
(${\geq}\SI{4.6}{\meter}$), so the long-band predictions and the long-band
hallucination guard are simply untestable on it. The downgrade from ``holds'' to
``partially holds'' is therefore a \emph{coverage gap in the held-out data, not a
weak result}: everything that could be tested passed, and the one thing that could
not be tested is a property of seq~00's car population, not of the method.

\section{Summary and limitations}
\label{sec:results-summary}

The completion method rests on three measured results. It adds genuinely unseen
surface --- about \SI{30}{\percent} of never-observed surface within
\SI{0.1}{\meter}, seven times the symmetry-mirror baseline, with essentially no
hallucination (Section~\ref{sec:results-coverage}). It improves the amodal box on
width, height, centre, and overall overlap, most where the sensor saw least, and
the length-geometry work turns an untreated length regression into a net gain
(Section~\ref{sec:results-box}). And both headline claims survive a
pre-registered replication on a held-out sequence
(Section~\ref{sec:results-heldout}). Moving cars, which the static-car metrics
cannot score, complete no worse than static ones
(Section~\ref{sec:results-movers}).

The limitations are stated as findings, each already quantified. The
\emph{far end} of the car is the weakest-covered region, and on long cars the
residual under-extension is a characterised, length-dependent limit that a
pre-registered fix could not remove without regressing compacts (Finding~\#45).
The per-car length estimate falls back to the fixed population prior on
\SI{23}{\percent} of final tracks --- those with fewer than five gate-passed
frames (Finding~\#41) --- which is above a comfortable threshold and is not
solved. The sparsest inputs complete more coarsely, and the axis-aligned plausibility
measure inflates the \emph{reported} width of diagonally-oriented cars
(Finding~\#44). And the two quantitative metrics measure static cars directly;
moving cars are covered only by the weaker plausibility check. None of these
overturns the headline --- completion adds unseen surface and a better box, and
generalises --- but they mark exactly where the method is soft, and they set the
agenda for the discussion that follows.
